% Propietario conservado: /Users/ruben/Documents/ChatGPT/jueces y controles/REVISION_ARGUMENTAL_APERTURAS_CIERRES_20260915/IV/ES/sections/iv_apendice_enumeraciones.tex
% RESULTADO_RECUPERADO. Fuentes conservadas en la entrega:
% antecedentes/controles_articulo_I/FIRMAS_EMISION.json:
% plus_classes, times_classes, seed_id_formula y coordinate_convention.
% antecedentes/controles_articulo_I/procedencia_excepcional.md:
% sección «Comprobación finita ejecutada», matriz A y seis generadores.
% Se explicitan las enumeraciones ya utilizadas; no se cambia su dominio.
\section{Finite enumerations of emissions and Witt stars}
\label{iv:iv:apendice-enumeraciones}

The two enumerations in this appendix concern distinct objects.
The first classifies the cursor signatures of APP emission;
the second computes the group generated by the stars of the Witt design.
Their domains and rules are retained so that each table can be reproduced
without being used as an input to the calculation.

\subsection{Initial-condition indices and signature recurrence}
\label{iv:iv:apendice-firmas}

Let \(i,j\in\{0,\ldots,8\}\) and \(a\in\{0,1,2,3\}\), with
direction order \((N,E,S,O)\) and vectors
\[
 v_N=(-1,0),\quad v_E=(0,1),\quad
 v_S=(1,0),\quad v_O=(0,-1).
\]
The index
\[
 \operatorname{id}(i,j,a)=4(9i+j)+a
\]
ranges over exactly \(\{0,\ldots,323\}\). Its inverse is obtained
from \(a=\operatorname{id}\bmod4\),
\(j=\lfloor\operatorname{id}/4\rfloor\bmod9\) and
\(i=\lfloor\operatorname{id}/36\rfloor\).
Two indices, one per sheet, enumerate the \(324^2\) joint conditions
in \eqref{iv:eq:semillas}; this domain is not identified with the 81 positions.

For each index, the cursor starts at \((i,j)\), with vector \(v_a\).
The instants \(t=1,\ldots,54\) are traversed in six windows
of nine instants. Write \(r_t=1+((t-1)\bmod9)\). The rule is:
\begin{enumerate}
\item On the additive sheet, an operation occurs when \(r_t\in\{1,4,7\}\):
\(\rho_9(i+j+2)\) is added to the window accumulator, and the
cursor moves by \(v_a\), modulo nine.
\item On the multiplicative sheet, an operation occurs when \(r_t\in\{2,5,8\}\):
\(\ell_9(\rho_9((i+1)(j+1)))\) is added, and the cursor moves
by \(v_a\), modulo nine. The extension of \(\ell_9\) by zero
outside the units, stated in Section~\ref{iv:sec:extension}, is retained.
\item At all other instants, that cursor does not move. After
instants \(27\) and \(54\), \(v_a\) is replaced by \(-v_a\).
\item At the end of each window, its accumulator is recorded modulo ten,
and only the accumulator is reset; neither position nor orientation is reset.
\end{enumerate}
This defines two functions on the 324 indices,
\[
 s:\{0,\ldots,323\}\longrightarrow\{0,\ldots,9\}^6,\qquad
 e:\{0,\ldots,323\}\longrightarrow\{0,\ldots,9\}^6.
\]
The multiplicity of a signature is the number of indices with that image.
The witness index in each table is the least index in its fibre.
Complete preimages are recovered by applying the recurrence to all
324 indices and retaining those that produce the stated signature.

\begin{longtable}{rcrr}
\caption{Additive signatures: image of \(s\), multiplicity and witness index.}
\label{iv:iv:tabla-firmas-aditivas}\\
\toprule
Class & Signature & Multiplicity & Witness index\\
\midrule
\endfirsthead
\toprule
Class & Signature & Multiplicity & Witness index\\
\midrule
\endhead
1&122564&18&17\\
2&122898&18&24\\
3&212645&18&5\\
4&212988&18&0\\
5&221456&18&29\\
6&221889&18&12\\
7&456221&18&28\\
8&465898&18&21\\
9&546988&18&9\\
10&564122&18&16\\
11&645212&18&4\\
12&654889&18&33\\
13&889221&18&13\\
14&889654&18&32\\
15&898122&18&25\\
16&898465&18&20\\
17&988212&18&1\\
18&988546&18&8\\
\bottomrule
\end{longtable}

\begin{longtable}{rcrr}
\caption{Multiplicative signatures: image of \(e\), multiplicity and witness index.}
\label{iv:iv:tabla-firmas-multiplicativas}\\
\toprule
Class & Signature & Multiplicity & Witness index\\
\midrule
\endfirsthead
\toprule
Class & Signature & Multiplicity & Witness index\\
\midrule
\endhead
1&000000&108&8\\
2&177393&8&1\\
3&177555&8&54\\
4&177771&8&11\\
5&339555&8&6\\
6&339771&8&15\\
7&339933&8&59\\
8&393177&8&0\\
9&393393&8&45\\
10&393555&8&40\\
11&555177&8&52\\
12&555339&8&4\\
13&555393&8&41\\
14&555555&24&84\\
15&555717&8&14\\
16&555771&8&18\\
17&555933&8&24\\
18&717555&8&12\\
19&717717&8&21\\
20&717933&8&19\\
21&771177&8&9\\
22&771339&8&13\\
23&771555&8&16\\
24&933339&8&57\\
25&933555&8&26\\
26&933717&8&17\\
\bottomrule
\end{longtable}

The sums of multiplicities are \(18\cdot18=324\) and
\(24\cdot8+24+108=324\). Joint emission is obtained by applying
\eqref{iv:eq:acoplamiento} to every pair of signatures. The inverse through
hundreds and tens proved in Section~\ref{iv:sec:extension} identifies the 468
emissions and gives their multiplicities as the product of the two fibres.
The subsequent map \eqref{iv:eq:psi-catalogo} produces the 243
visible words. The initial domain, emission, its fibres and the ternary
quotient are thus specified separately.

\subsection{Explicit generators and closure of the stars}
\label{iv:iv:apendice-estrellas}

The matrix \(A_W\) of \eqref{iv:exc:aw} is retained, and the following words
are formed in lexicographic order:
\[
 c(w)=(w,wA_W),\qquad w\in\mathbb F_3^6.
\]
They are grouped by weight and the supports of weight six are retained.
This produces the enumerator \eqref{iv:exc:enumerador} and the 132 hexads
of \(\mathcal H\). For each quadruple \(B\subset\{1,\ldots,12\}\),
the four pairs \(H\setminus B\), with \(B\subset H\in\mathcal H\),
define \(\sigma_B\): it fixes \(B\) and interchanges the two points
of each pair. This gives the family of 495 stars of
proposition~\ref{iv:exc:m12}.

The six generators used in its closure are expressed by their lists
of images of \(1,\ldots,12\):
\[
\begin{array}{c|rrrrrrrrrrrr}
 &1&2&3&4&5&6&7&8&9&10&11&12\\ \hline
g_1&1&2&3&4&12&11&9&10&7&8&6&5\\
g_2&1&2&3&12&5&10&11&9&8&6&7&4\\
g_3&1&2&3&11&10&6&8&7&12&5&4&9\\
g_4&1&2&12&4&5&9&8&7&6&11&10&3\\
g_5&1&12&3&4&5&8&10&6&11&7&9&2\\
g_6&12&2&3&4&5&7&6&11&10&9&8&1
\end{array}
\]
Each row is verified to be a star by comparison with the
\(\sigma_B\) already constructed. For \(j=1,\ldots,6\), the closure
is computed by
\[
 G_j^{(0)}=\{\operatorname{id}\},\qquad
 G_j^{(n+1)}
 =G_j^{(n)}\cup
   \{g_i\circ p:1\leq i\leq j,\ p\in G_j^{(n)}\}.
\]
The increasing sequence stabilizes in the finite set
\(\operatorname{Sym}(12)\). Since the generators are involutions,
its stable value is exactly the group \(G_j\) they generate.
With the convention \((g_i\circ p)(a)=g_i(p(a))\), their orders are
\[
\begin{array}{c|rrrrrr}
 j&1&2&3&4&5&6\\ \hline
 |G_j|&2&6&18&360&7920&95040 .
\end{array}
\]
Finally, one checks that each of the 495 stars belongs
to \(G_6\) and that they all preserve \(\mathcal H\). The first inclusion,
together with membership of the six generators in the family,
makes the two generated groups equal. Preservation is checked
through \(\{\sigma_B(H):H\in\mathcal H\}=\mathcal H\).
This is the finite enumeration used in the proof of
\ref{iv:exc:m12}; subsequent recognition as \(M_{12}\)
retains the external reference declared there.

